\documentclass[12pt]{article}

\usepackage[a4paper, margin=1in]{geometry}
\usepackage{amsmath, amssymb, amsthm}
\usepackage{setspace}
\usepackage{titlesec}
\usepackage{enumitem}
\usepackage{hyperref}
\usepackage{fancyhdr}

\setstretch{1.15}
\setlength{\parskip}{0.5em}
\setlength{\parindent}{0pt}

\pagestyle{fancy}
\fancyhf{}
\rhead{CSE400}
\lhead{Lecture 21 \& 22}
\cfoot{\thepage}

\titleformat{\section}{\large\bfseries}{\thesection.}{0.5em}{}
\titleformat{\subsection}{\normalsize\bfseries}{\thesubsection.}{0.5em}{}

\title{\textbf{CSE400 -- Fundamentals of Probability in Computing}\\
\Large Lecture 21 \& 22 Scribe\\
\large Inequalities and Tail Bounds}
\author{}
\date{}

\begin{document}

\maketitle
\vspace{-1.5em}
\hrule
\vspace{1em}

\section{Introduction}

In probability and computing systems, it is often important to understand not only the \textbf{average behavior} of a random variable, but also the probability of \textbf{rare and extreme events}. Such events are captured through the \textbf{tail} of a distribution. These lectures focus on the notion of tails, their importance in computing systems, and standard inequalities used to bound tail probabilities when exact computation is difficult.

\section{Tail of a Random Variable}

\subsection{Definition}
The \textbf{tail} of a random variable \(X\) is defined as
\[
P(X > x)
\]
This quantity measures the probability that \(X\) exceeds a given threshold and lies in the extreme right side of the distribution.

\section{Motivation for Studying Tails}

Tail probabilities are important because they correspond to \textbf{rare but critical events} in real systems. Examples discussed in the lecture include:
\begin{itemize}[leftmargin=1.5em]
    \item Jobs delayed more than 24 hours
    \item Packets dropped due to full buffer
\end{itemize}

The lecture emphasizes the following distinction:
\begin{itemize}[leftmargin=1.5em]
    \item The \textbf{mean} describes average behavior and is generally easier to compute.
    \item The \textbf{tail} captures rare extreme behavior and is usually harder to compute.
\end{itemize}

The main difficulty arises because exact tail probabilities often require summing a large number of very small probabilities.

\section{Examples of Tail Events}

\subsection{Jobs Distributed Across Servers}
Suppose \(n\) jobs are distributed randomly among \(n\) servers. Let \(X\) denote the number of jobs assigned to a particular server.

Then,
\[
X \sim \text{Binomial}(n,p)
\]

The probability that this server gets at least \(k\) jobs is
\[
P(X \ge k)=\sum_{i=k}^{n}\binom{n}{i}p^i(1-p)^{n-i}
\]

This is a \textbf{tail event}, representing a rare overload scenario.

\subsection{Poisson Arrivals at a Datacenter}
Suppose jobs arrive to a datacenter according to a Poisson process with rate \(\lambda\) jobs per hour. Let \(X\) denote the number of arrivals in one hour.

Then,
\[
X \sim \text{Poisson}(\lambda)
\]

The probability of receiving at least \(k\) arrivals is
\[
P(X \ge k)=\sum_{i=k}^{\infty} e^{-\lambda}\frac{\lambda^i}{i!}
\]

This is also a \textbf{tail event}, corresponding to a burst in traffic or load.

\section{Why Tail Bounds Are Needed}

Exact tail computation is often difficult due to complicated summations such as
\[
P(X \ge k)=\sum_{i=k}^{\infty} e^{-\lambda}\frac{\lambda^i}{i!}
\]
or
\[
P(X \ge k)=\sum_{i=k}^{n}\binom{n}{i}p^i(1-p)^{n-i}
\]

Hence, rather than computing exact values, the lecture proposes finding \textbf{upper bounds} of the form
\[
P(X \ge k)\le \text{something simple}
\]

These bounds are useful because they are:
\begin{itemize}[leftmargin=1.5em]
    \item easier to compute,
    \item guaranteed to be valid,
    \item and practically useful in system analysis.
\end{itemize}

\section{Running Example}

To compare different inequalities, the lecture uses the same running example throughout.

\subsection{Experiment}
Flip a fair coin \(n\) times.

\subsection{Random Variable}
\[
X = \text{number of heads}
\]

Then,
\[
X \sim \text{Binomial}\left(n,\frac12\right)
\]

\subsection{Mean}
\[
E[X]=\mu=\frac{n}{2}
\]

\subsection{Question}
We wish to bound
\[
P\left(X \ge \frac34 n\right)
\]

This corresponds to the probability of obtaining \textbf{unusually many heads}.

\section{Markov's Inequality}

\subsection{Statement}
For a \textbf{non-negative} random variable \(X \ge 0\) and any \(a>0\),
\[
P(X \ge a)\le \frac{E[X]}{a}
\]

This inequality provides an upper bound on the \textbf{right tail} using only the mean.

\subsection{Interpretation}
Markov’s inequality expresses the idea that if the average value of a non-negative random variable is small, then the probability of it taking a very large value must also be small.

\subsection{Proof}
Starting from the expectation:
\[
E[X]=\sum_{x=0}^{\infty} x\,p_X(x)
\]

Restricting to the values \(x \ge a\), we obtain
\[
E[X]\ge \sum_{x=a}^{\infty} x\,p_X(x)
\]

Since \(x \ge a\) in this range,
\[
E[X]\ge \sum_{x=a}^{\infty} a\,p_X(x)
\]

\[
E[X]\ge a\sum_{x=a}^{\infty} p_X(x)
\]

\[
E[X]\ge aP(X \ge a)
\]

Therefore,
\[
P(X \ge a)\le \frac{E[X]}{a}
\]

\subsection{Markov on the Running Example}
For
\[
X \sim \text{Binomial}\left(n,\frac12\right)
\]
we have
\[
E[X]=\frac{n}{2}
\]

We want
\[
P\left(X \ge \frac34 n\right)
\]

Applying Markov’s inequality:
\[
P\left(X \ge \frac34 n\right)\le \frac{\frac n2}{\frac34 n}
\]

\[
=\frac23
\]

\textbf{Observation:} This is a weak bound, since it does not improve as \(n\) increases.

\subsection{Example: Income Distribution}
Suppose the average income is
\[
E[X]=10{,}000
\]

To bound the fraction of people earning at least \(50{,}000\), apply Markov:
\[
P(X \ge 50{,}000)\le \frac{10{,}000}{50{,}000}
\]

\[
=0.2
\]

Hence, at most \(20\%\) of people can earn \(50{,}000\) or more.

\subsection{Solved Example 1: Exam Grades}
Let \(X\) denote the student grade. Suppose:
\[
E[X]=70
\]
and the ``A'' cutoff is \(90\).

Then
\[
P(X \ge 90)\le \frac{70}{90}
\]

\[
=\frac79\approx 0.778
\]

Thus, at most \(77.8\%\) of students can receive an ``A''.

\subsection{Solved Example 2: Server Request Load}
Suppose a server processes an average of \(5{,}000\) requests per second. Let \(X\) denote the number of requests received in one second.

Then
\[
E[X]=5000
\]

We want to bound
\[
P(X \ge 20{,}000)
\]

Applying Markov:
\[
P(X \ge 20{,}000)\le \frac{5000}{20{,}000}
\]

\[
=\frac14=0.25
\]

Thus, the probability is at most \(25\%\).

\subsection{Limitation of Markov's Inequality}
Markov’s inequality only provides bounds on the \textbf{right tail}:
\[
P(X \ge a)
\]

It does not directly bound probabilities such as
\[
P(X<a)
\]
without additional information.

\section{Chebyshev's Inequality}

\subsection{Statement}
For a random variable \(X\) with mean \(\mu\) and variance \(\operatorname{Var}(X)\), and for any \(a>0\),
\[
P(|X-\mu|\ge a)\le \frac{\operatorname{Var}(X)}{a^2}
\]

Equivalently, if \(\sigma^2 = \operatorname{Var}(X)\), then
\[
P(|X-\mu|\ge a)\le \frac{\sigma^2}{a^2}
\]

This inequality bounds the probability that \(X\) lies \textbf{far from its mean}.

\subsection{Interpretation}
Chebyshev’s inequality uses not only the mean, but also the \textbf{variance}, making it stronger than Markov in many settings.

\subsection{Proof Idea}
Apply Markov’s inequality to the non-negative random variable
\[
(X-\mu)^2
\]

Then,
\[
P(|X-\mu|\ge a)=P((X-\mu)^2\ge a^2)
\]

Applying Markov gives
\[
P((X-\mu)^2\ge a^2)\le \frac{E[(X-\mu)^2]}{a^2}
\]

Since
\[
E[(X-\mu)^2]=\operatorname{Var}(X)
\]

we obtain
\[
P(|X-\mu|\ge a)\le \frac{\operatorname{Var}(X)}{a^2}
\]

\subsection{Chebyshev on the Running Example}
For
\[
X \sim \text{Binomial}\left(n,\frac12\right)
\]
we have
\[
\mu=\frac n2,\qquad \operatorname{Var}(X)=\frac n4
\]

We wish to bound
\[
P\left(X \ge \frac34 n\right)
\]

From the lecture:
\[
P\left(X \ge \frac34 n\right)
=
\frac12 P\left(\left|X-\frac n2\right|\ge \frac n4\right)
\]

Applying Chebyshev:
\[
P\left(\left|X-\frac n2\right|\ge \frac n4\right)\le \frac{n/4}{(n/4)^2}
\]

\[
=\frac{4}{n}
\]

Thus,
\[
P\left(X \ge \frac34 n\right)\le \frac12\cdot \frac{4}{n}
\]

\[
=\frac{2}{n}
\]

\textbf{Observation:} This bound decreases with \(n\) and is therefore better than Markov’s bound.

\subsection{Example: Exam Scores}
Suppose exam scores have
\[
\mu=70,\qquad \sigma=10
\]

We want to bound the fraction of students scoring below \(40\) or above \(100\). This is
\[
P(|X-70|\ge 30)
\]

Applying Chebyshev:
\[
P(|X-70|\ge 30)\le \frac{10^2}{30^2}
\]

\[
=\frac{100}{900}=\frac19
\]

Hence, at most \(\frac19\) of students are this far from the mean.

\subsection{Solved Example: Exam Grades with Variance}
Suppose
\[
\mu=70,\qquad \sigma=5,\qquad \sigma^2=25
\]
and the ``A'' cutoff is \(90\).

Observe that
\[
X\ge 90 \implies |X-70|\ge 20
\]

Hence,
\[
P(X\ge 90)\le P(|X-70|\ge 20)
\]

Applying Chebyshev:
\[
P(|X-70|\ge 20)\le \frac{25}{20^2}
\]

\[
=\frac{25}{400}=\frac{1}{16}=0.0625
\]

Thus,
\[
P(X\ge 90)\le 0.0625
\]

This is a much tighter bound than the Markov bound of \(0.778\).

\section{Chernoff Bound and Poisson Tail}

The lecture outline explicitly includes:
\begin{itemize}[leftmargin=1.5em]
    \item Chernoff Bound
    \item Poisson Tail
    \item Definition and proof
    \item Example
\end{itemize}

\subsection{Poisson Tail}
For
\[
X \sim \text{Poisson}(\lambda)
\]
the upper tail is
\[
P(X \ge k)=\sum_{i=k}^{\infty} e^{-\lambda}\frac{\lambda^i}{i!}
\]

This expression is used to quantify rare arrival bursts in queueing and server systems.

\subsection{Role of Chernoff Bound}
The lecture presents Chernoff bound as a \textbf{progressively tighter tail bound} than both Markov and Chebyshev.

Thus, the sequence of inequalities emphasized in the lecture is:
\[
\text{Markov} \rightarrow \text{Chebyshev} \rightarrow \text{Chernoff}
\]

where each step gives a better tail estimate.

\section{Jensen's Inequality}

The lecture outline includes:
\begin{itemize}[leftmargin=1.5em]
    \item Jensen’s Inequality
    \item Convex Functions
    \item Proof and Example
\end{itemize}

\subsection{Convex Functions}
Jensen’s inequality is introduced in the context of \textbf{convex functions}.

\subsection{Role in These Lectures}
Jensen’s inequality is included as an important inequality involving:
\begin{itemize}[leftmargin=1.5em]
    \item random variables,
    \item expectation,
    \item and convex functions.
\end{itemize}

\section{Case Study and System-Level Understanding}

The lecture concludes with a case study centered around the question:

\subsection*{How to predict the worst-case demands on server?}

This highlights the practical role of tail bounds in computing systems. Tail probabilities are essential in estimating rare but important events such as:
\begin{itemize}[leftmargin=1.5em]
    \item overloads,
    \item bursts of requests,
    \item queue buildup,
    \item and performance degradation.
\end{itemize}

\textbf{Main Insight:} Average behavior alone is insufficient for system design. One must also understand the probability of rare extreme events.

\section{Final Summary}

\subsection{Core Definitions}
Tail of \(X\):
\[
P(X>x)
\]

Right-tail event:
\[
P(X\ge a)
\]

Deviation from mean:
\[
P(|X-\mu|\ge a)
\]

\subsection{Exact Tail Forms}
\textbf{Binomial:}
\[
P(X \ge k)=\sum_{i=k}^{n}\binom{n}{i}p^i(1-p)^{n-i}
\]

\textbf{Poisson:}
\[
P(X \ge k)=\sum_{i=k}^{\infty} e^{-\lambda}\frac{\lambda^i}{i!}
\]

\subsection{Markov's Inequality}
\[
P(X\ge a)\le \frac{E[X]}{a}
\qquad (X\ge 0)
\]

\subsection{Chebyshev's Inequality}
\[
P(|X-\mu|\ge a)\le \frac{\operatorname{Var}(X)}{a^2}
\]

\subsection{Running Example}
\[
X\sim \text{Binomial}\left(n,\frac12\right),\quad E[X]=\frac n2,\quad \operatorname{Var}(X)=\frac n4
\]

Tail of interest:
\[
P\left(X\ge \frac34 n\right)
\]

Markov bound:
\[
P\left(X\ge \frac34 n\right)\le \frac23
\]

Chebyshev bound:
\[
P\left(X\ge \frac34 n\right)\le \frac2n
\]

\subsection{Final Conceptual Message}
\begin{itemize}[leftmargin=1.5em]
    \item Exact tail probabilities are often difficult to compute.
    \item Tail inequalities provide simpler, guaranteed bounds.
    \item These bounds are important for understanding \textbf{rare extreme events} in computing systems.
\end{itemize}

\section{One-Page Exam Memory Sheet}

\[
P(X>x)
\]

\[
P(X\ge a)\le \frac{E[X]}{a}
\]

\[
P(|X-\mu|\ge a)\le \frac{\operatorname{Var}(X)}{a^2}
\]

\[
P(X\ge k)=\sum_{i=k}^{n}\binom{n}{i}p^i(1-p)^{n-i}
\]

\[
P(X\ge k)=\sum_{i=k}^{\infty} e^{-\lambda}\frac{\lambda^i}{i!}
\]

\[
X\sim \text{Binomial}\left(n,\frac12\right)
\]

\[
E[X]=\frac n2,\qquad \operatorname{Var}(X)=\frac n4
\]

\[
P\left(X\ge \frac34 n\right)\le \frac23
\]

\[
P\left(X\ge \frac34 n\right)\le \frac2n
\]

\[
\text{Markov} \rightarrow \text{Chebyshev} \rightarrow \text{Chernoff}
\]

\end{document}